\subsection{Individual Expression Under Image and Initialization Conditions}
\label{sec:individual}

Annotators differ in both retained structure and boundary position. The fireplace example contrasts an eight-turn recess with a four-turn enclosure (Figure~\ref{fig:individual-fireplace}); even annotators retaining the same turns place their endpoints differently. The out-of-building profiles distinguish simplifying (C1), typical (C2), and turn-retaining (C3) histories (Figure~\ref{fig:continuous-profiles}).

\begin{figure*}[t]
\centering
\setcounter{subfigure}{0}
\begin{minipage}{150mm}
\centering
\ESPanel{.323\linewidth}{panels/fig09/a.pdf}{Structural response}%
\hfill%
\ESPanel{.323\linewidth}{panels/fig09/b.pdf}{Quality ranks}%
\hfill%
\ESPanel{.323\linewidth}{panels/fig09/c.pdf}{Difficult panorama}

\par\vspace{3pt}
\ESPanel{.323\linewidth}{panels/fig09/d.pdf}{C1: A55}%
\hfill%
\ESPanel{.323\linewidth}{panels/fig09/e.pdf}{C2: A44}%
\hfill%
\ESPanel{.323\linewidth}{panels/fig09/f.pdf}{C3: A43}

\end{minipage}
\caption{Structural response, relative quality and original drawings. (a--b) Annotator-level structural responses and quality ranks across difficulty conditions, using 58 and 59 people respectively; each dot represents one person. Relative Q is centered on the image--target peer median; higher within-image ranks indicate better relative quality. (c) The difficult panorama. (d--f) Submissions from C1--C3 on the panorama in (c) retain five, six and eight turns on a common BEV scale; gray dashes mark the reference and thick segments mark short walls.}
\label{fig:continuous-profiles}
\label{fig:class-annotations}
\end{figure*}


Differences are larger on images with short walls or self-occlusion. Across 33 six-turn reference panoramas (1,023 submissions), image-weighted simplification rates were 10.7\% without short edges or self-occlusion, 29.8\% with short edges, and 55.7\% under self-occlusion. C3 annotators increased their marked turns from easy to difficult scenes by 2.330 more turns than C2 annotators, with the difference positive in all 15 buildings. Their advantage also involved localization: mean matched-boundary error was $0.941^\circ$ for C3 and $1.730^\circ$ for C2 on 201 shared images. Across 205 images, C3 exceeded C2 by 14.178 $Q$ points on average, with gains on 188 images, including easy scenes where their turn counts were similar. Structural habits therefore describe both responses to geometric difficulty and differences in the resulting geometry. Joint profiles also distinguish recorded effort: J2 has a mean $Q$ difference of $-0.133$ and a geometric mean duration ratio of 0.885 relative to J1 across 162 shared images. Supplement D1 and Supplements C7 and C19 report the extended comparisons.

Initialized editing produces a different set of outcomes. Three participants submitted the natural initialization for \texttt{e9zR4mvMWw7-39} unchanged, retaining errors identified before the study. The six controlled-defect tasks yielded 156 opportunities from 26 annotators, including 138 standard primary-target submissions: 20 were unchanged, 32 precisely removed an injected duplicate pair, one edited other coordinates but retained that pair, and 85 edited coordinates without cleanly resolving the target defect (Supplements B8, B20 and C11). Editing activity and successful correction are distinct outcomes.

Task reports also relate differently to the two conditions. On 72 Manual images, participants reporting occlusion omitted 4.81 percentage points less reference floor area than non-reporters on the same image (95\% building-bootstrap CI $[-7.47,-1.24]$); the direction persisted in within-person comparisons. Semi reports showed no consistent association with floor omission (Supplement C11). These records connect reported visual difficulty to submitted geometry while distinguishing drawing from initialized editing.

